\section{Adjoint boundaries from star-compatible coalgebras}
\label{sec:glm_coalgebra_adjoint}

Let $H$ be a complex coalgebra with a conjugate-linear involution such that
$\Delta(x^*)=(*\otimes*)\Delta(x)$. This is the coproduct compatibility
in the definition of a $C^*$-weak Hopf algebra in
\cite[Definition~5.8]{Molnar2022MPOAlgebrasI}.
The construction below uses only this coalgebra, a physical linear map,
and a dual representation; it requires neither an antipode nor an integral.
The involution $\kappa$ below preserves convolution order and differs from
the dual $C^*$-algebra involution, which also involves the antipode
\cite[Equation~(5.38)]{Molnar2022MPOAlgebrasI}.

\begin{theorem}[Same-order involution of the convolution dual]
    \label{thm:glm_coalgebra_dual_star}
    \lean{Coalgebra.dualStar_mul,Coalgebra.dualStarMulEquiv,
        Coalgebra.dualStar_one}
    \leanok
    For $f\in H^*$, define $\kappa(f)(x)=\overline{f(x^*)}$.
    Then $\kappa$ is conjugate-linear, involutive, and satisfies
    \begin{align}
        \kappa(fg) & =\kappa(f)\kappa(g), \notag \\
        \kappa(1)  & =1.
        \label{eq:glm_dual_star_product}
    \end{align}
    Here multiplication is convolution, and $1$ is the counit.
\end{theorem}
\begin{proof}
    \leanok
    On a pure tensor, complex conjugation sends
    $f(x^*)g(y^*)$ to $\kappa(f)(x)\kappa(g)(y)$.
    Extend this identity additively to $\Delta(x)$ and use the
    coproduct compatibility. Involutivity makes $\kappa$ a bijective
    multiplicative map, so it fixes the unit.
\end{proof}

\begin{definition}[Tensor from physical coefficients]
    \label{def:glm_dual_representation_tensor}
    \lean{MPOTensor.dualCoefficient,MPOTensor.ofDualRepresentation}
    \leanok
    Let $\phi:H\to M_d(\mathbb C)$ be linear and let
    $\psi:H^*\to M_D(\mathbb C)$ be a unital algebra representation.
    Define $f_{ij}(x)=\phi(x)_{ij}$ and $T^{ij}=\psi(f_{ij})$.
\end{definition}

\begin{theorem}[Ordered coefficient products]
    \label{thm:glm_dual_representation_words}
    \lean{MPOTensor.evalWord_ofDualRepresentation,
        MPOTensor.star_dualCoefficient,MPOTensor.star_dualCoefficient_prod}
    \leanok
    \uses{def:glm_dual_representation_tensor}
    Suppose $\phi(x^*)=\phi(x)^\dagger$. For every pair of physical
    configurations $\sigma,\tau$ of length $N\geq0$, put
    $f_{\sigma,\tau}=f_{\sigma_0\tau_0}\cdots f_{\sigma_{N-1}\tau_{N-1}}$,
    with the empty product equal to $1$. Then
    \begin{align}
        T^{\sigma_0\tau_0}\cdots T^{\sigma_{N-1}\tau_{N-1}}
                                & =\psi(f_{\sigma,\tau}), \notag \\
        \kappa(f_{\sigma,\tau}) & =f_{\tau,\sigma}.
        \label{eq:glm_dual_representation_words}
    \end{align}
\end{theorem}
\begin{proof}
    \leanok
    \uses{thm:glm_coalgebra_dual_star}
    Star compatibility gives $\kappa(f_{ij})=f_{ji}$.
    Induction on $N$, using multiplicativity and unitality of $\psi$
    and (\ref{eq:glm_dual_star_product}), proves both identities.
\end{proof}

In the following diagrams, a west virtual leg is a matrix row and an
east virtual leg is a matrix column. The north physical leg $\sigma_n$
is an output and the south physical leg $\tau_n$ is an input, so the
box at site $n$ is the matrix $T^{\sigma_n\tau_n}$. The first identity
of (\ref{eq:glm_dual_representation_words}) reads as follows: the open
chain of tensor letters is one matrix of the representation, evaluated
at the ordered convolution product of the coefficient functionals.
% TENKZ-COALGEBRA-WORDS-BEGIN
% Formula: (T^{sigma_0 tau_0} ... T^{sigma_{N-1} tau_{N-1}})_{alpha beta}
% = psi(f_{sigma_0 tau_0} ... f_{sigma_{N-1} tau_{N-1}})_{alpha beta},
% the first identity of eq:glm_dual_representation_words, coefficientwise
% in fixed configurations sigma, tau of length N.
% Ink-to-index map: box n is T^{sigma_n tau_n}; its north physical leg
% sigma_n is the output (row) index and its south physical leg tau_n the
% input (column) index of phi(x). West ports are matrix rows, east ports
% are matrix columns. External alpha, beta range over Fin D.
% Contracted legs: the D-dimensional virtual bonds between adjacent boxes.
% Three sites are drawn and the dots atom stands for the remaining ones.
% On the right, sigma and tau are labels of the functional, not legs.
% Boundary signature: left open:w, open:e, phys:n=N, phys:s=N; right
% open:w, open:e.
\begin{tenkzequation}
    \begin{tenkz}[rows={wire}, cols=4, bonds=none]
        \tn[at=(1,1),name=T0,skin=box,size=l,
            ports={180:virtual:$\alpha$,0:virtual,
                    90:physical:$\sigma_0$,270:physical:$\tau_0$}]{T}
        \tn[at=(1,2),name=T1,skin=box,size=l,
            ports={180:virtual,0:virtual,
                    90:physical:$\sigma_1$,270:physical:$\tau_1$}]{T}
        \tn[at=(1,3),name=D,skin=dots,
            ports={180:virtual,0:virtual}]{}
        \tn[at=(1,4),name=TN,skin=box,size=l,
            ports={180:virtual,0:virtual:$\beta$,
                    90:physical:$\sigma_{N-1}$,
                    270:physical:$\tau_{N-1}$}]{T}
        \tnwire{T0.0}{T1.180}
        \tnwire{T1.0}{D.180}
        \tnwire{D.0}{TN.180}
    \end{tenkz}
    $=$
    \begin{tenkz}[rows={wire}, cols=2, bonds=none]
        \tn[at=(1,1),name=P,skin=box,size=l,wide=2,
            ports={180:virtual:$\alpha$,0:virtual:$\beta$}]{\psi(f_{\sigma,\tau})}
    \end{tenkz}
\end{tenkzequation}
% TENKZ-COALGEBRA-WORDS-END

\begin{theorem}[One adjoint boundary at every length]
    \label{thm:glm_coalgebra_adjoint_boundary}
    \lean{MPOTensor.exists_boundary_trace_dualStar,
        MPOTensor.exists_adjointBoundary_of_dualRepresentation,
        MPOTensor.isBoundaryAdjointClosed_of_dualRepresentation}
    \leanok
    \uses{def:glm_dual_representation_tensor,def:glm_boundary_adjoint_closed}
    Suppose $\phi$ preserves the involution and $\psi$ is injective.
    For every $X\in M_D(\mathbb C)$ there exists one
    $Y\in M_D(\mathbb C)$ such that
    \begin{align}
        (O^N_{T,X})^\dagger & =O^N_{T,Y}
        \label{eq:glm_coalgebra_adjoint_boundary}
    \end{align}
    for every $N\geq0$. In particular, $T$ is boundary adjoint closed.
\end{theorem}
\begin{proof}
    \leanok
    \uses{thm:glm_dual_representation_words}
    The functional
    $\ell_X(f)=\overline{\operatorname{tr}(X\psi(\kappa(f)))}$
    is complex linear. Injectivity of $\psi$ allows it to extend to a
    linear functional $\theta$ on $M_D(\mathbb C)$. The nondegenerate
    trace pairing provides $Y$ with $\theta(M)=\operatorname{tr}(YM)$.
    Consequently
    \begin{align*}
        \operatorname{tr}(Y\psi(f_{\sigma,\tau}))
         & =\overline{\operatorname{tr}
        (X\psi(\kappa(f_{\sigma,\tau})))}
        =\overline{\operatorname{tr}(X\psi(f_{\tau,\sigma}))}.
    \end{align*}
    This proves (\ref{eq:glm_coalgebra_adjoint_boundary}) entrywise.
    The boundary $Y$ was chosen before $N$. At $N=0$, the identity is
    $\operatorname{tr}(Y)=\overline{\operatorname{tr}(X)}$.
\end{proof}

Coefficientwise, (\ref{eq:glm_coalgebra_adjoint_boundary}) says
$\operatorname{tr}(YT^{\sigma_0\tau_0}\cdots T^{\sigma_{N-1}\tau_{N-1}})
    =\overline{\operatorname{tr}(XT^{\tau_0\sigma_0}\cdots
        T^{\tau_{N-1}\sigma_{N-1}})}$.
Conjugating every factor of the right-hand trace gives the following
picture for $N\geq1$: the adjoint exchanges the north and south
physical legs of every letter and conjugates every entry, with no
virtual transpose, and the result is the same ring closed by the
single boundary $Y$.
% TENKZ-COALGEBRA-ADJOINT-BEGIN
% Formula: tr(Y T^{sigma_0 tau_0} ... T^{sigma_{N-1} tau_{N-1}})
% = tr(conj(X) conj(T^{tau_0 sigma_0}) ... conj(T^{tau_{N-1} sigma_{N-1}})),
% the (sigma, tau) coefficient of eq:glm_coalgebra_adjoint_boundary,
% where conj is entrywise complex conjugation.
% Ink-to-index map: a box M has west virtual row index and east virtual
% column index; adjacent east/west legs are contracted, and the periodic
% closure contracts the final east leg with the first west leg, so the
% row reads as a trace from left to right. North physical legs are
% output indices and south physical legs are input indices: on the left,
% site n carries sigma_n north and tau_n south; on the right the same
% labels appear exchanged, tau_n north and sigma_n south, on the
% conjugated letters. The boundary boxes have virtual legs only.
% Three sites are drawn and the dots atom stands for the remaining ones.
% Boundary signature on both sides: virtual-west=0, virtual-east=0,
% physical-north=N, physical-south=N.
\begin{tenkzequation}
    \begin{tenkz}[rows={op}, cols=5, boundary=periodic]
        \tn[skin=box,ports={180:virtual,0:virtual}]{Y} &
        \tn[skin=mpo,ports={180:virtual,0:virtual,
                    90:physical:$\sigma_0$,270:physical:$\tau_0$}]{T} &
        \tn[skin=mpo,ports={180:virtual,0:virtual,
                    90:physical:$\sigma_1$,270:physical:$\tau_1$}]{T} &
        \tn[skin=dots,ports={180:virtual,0:virtual}]{} &
        \tn[skin=mpo,ports={180:virtual,0:virtual,
                    90:physical:$\sigma_{N-1}$,
                    270:physical:$\tau_{N-1}$}]{T}
    \end{tenkz}
    $=$
    \begin{tenkz}[rows={op}, cols=5, boundary=periodic]
        \tn[skin=box,ports={180:virtual,0:virtual}]{\overline X} &
        \tn[skin=mpo,ports={180:virtual,0:virtual,
                    90:physical:$\tau_0$,270:physical:$\sigma_0$}]{\overline T} &
        \tn[skin=mpo,ports={180:virtual,0:virtual,
                    90:physical:$\tau_1$,270:physical:$\sigma_1$}]{\overline T} &
        \tn[skin=dots,ports={180:virtual,0:virtual}]{} &
        \tn[skin=mpo,ports={180:virtual,0:virtual,
                    90:physical:$\tau_{N-1}$,
                    270:physical:$\sigma_{N-1}$}]{\overline T}
    \end{tenkz}
\end{tenkzequation}
% TENKZ-COALGEBRA-ADJOINT-END

\begin{theorem}[Canonical parent symmetry of a dual representation]
    \label{thm:glm_dual_representation_parent}
    \lean{MPOTensor.IsBoundaryCompatible.parentInteractionES_commute_of_dualRepresentation,
        MPOTensor.IsBoundaryCompatible.parentHamiltonianES_commute_of_dualRepresentation}
    \leanok
    \uses{def:glm_dual_representation_tensor,def:glm_boundary_closedness,
        def:parent_interaction,def:parent_hamiltonian,def:mpdo_commuting_boundary}
    Under the hypotheses of Theorem~\ref{thm:glm_coalgebra_adjoint_boundary},
    suppose $T$ is arbitrary-boundary compatible with a state tensor $A$.
    For $L>0$, the canonical local parent commutes with $O^L_{T,X}$
    for every $X$. For $N\geq L$, the periodic canonical parent
    commutes with $O^N_{T,X}$ whenever $X$ commutes with every tensor
    letter $T^{ij}$.
\end{theorem}
\begin{proof}
    \leanok
    \uses{thm:glm_coalgebra_adjoint_boundary,thm:glm_boundary_parent_local,
        thm:glm_boundary_parent_periodic}
    Compatibility makes the local state support invariant. Its
    orthogonal complement is invariant as well, since the entire boundary
    family is adjoint closed. The complementary orthogonal projector
    therefore satisfies $[h_L,O^L_{T,X}]=0$ by
    Theorem~\ref{thm:glm_boundary_parent_local}. The
    tensor-letter commutation condition transports this equality through
    cyclic windows, including those crossing the periodic boundary;
    Theorem~\ref{thm:glm_boundary_parent_periodic} then gives
    $[H_{L,N},O^N_{T,X}]=0$.
\end{proof}
